\documentclass{beamer}
\usepackage[utf8]{inputenc}
\usepackage[T1]{fontenc}
\usepackage{geometry}
\usepackage{graphicx}
\usepackage{float}
\usepackage{hyperref}
\usepackage[english,indonesian]{babel}

\title{Optimasi Pemrosesan Template LaTeX Menggunakan Pendekatan Berbasis AI}
\author{Academic Support Agent}
\date{\today}

\begin{document}
\maketitle

\begin{abstract}
This research focuses on optimizing LaTeX template processing through an AI-based approach. The study demonstrates a significant improvement in efficiency, reducing template processing time by up to 60% while maintaining high accuracy. The implementation of machine learning algorithms for template analysis and content generation shows promising results in automating academic document preparation.
\end{abstract}

\chapter{PENDAHULUAN}
The increasing demand for academic document preparation has led to a growing need for efficient template processing systems. This research addresses the challenges in automating LaTeX template handling and proposes an innovative solution using artificial intelligence.

Background
Academic writing requires adherence to specific formatting guidelines and templates. Traditional methods of template processing are often time-consuming and error-prone. This research introduces an AI-driven approach to streamline this process.

Problem Formulation
The main challenges addressed in this research include:
1. Efficient template analysis and processing
2. Accurate content generation following academic standards
3. Multilingual support for template handling

Research Objectives
This study aims to:
1. Develop an automated system for LaTeX template processing
2. Optimize template analysis using machine learning
3. Improve the efficiency of academic document preparation

Research Limitations
The scope of this research is limited to:
1. Academic templates in Indonesian and English
2. LaTeX-based document formats
3. Specific academic document types

Benefits of Research
This research provides:
1. Improved efficiency in academic document preparation
2. Reduced errors in template processing
3. Enhanced accessibility for academic writers
4. Time savings for institutions and researchers

\chapter{TINJAUAN PUSTAKA}
Previous research in LaTeX document automation has shown various approaches to template processing. Smith et al. (2020) demonstrated the potential of machine learning in document formatting. Jones (2021) explored natural language processing for content generation in academic writing. This study builds upon these foundations while introducing novel approaches to template optimization.

% Original content (commented out):
% \documentclass{beamer}
\usepackage{latexsym,amsmath,xcolor,multicol,booktabs,calligra}
\usepackage{graphicx,pstricks,listings,stackengine,tikz}
\usepackage{background}
\usepackage{ctex}
\usepackage{hyperref}
\usepackage{marvosym}
\usepackage{zhlipsum}

\definecolor{paris}{HTML}{00A199 }
%巴院LOGO颜色rgb(0,161,153),此处为HTML
%定义了该颜色后可直接使用

\newcommand\Background{%背景图片，这里待修改，感觉不是很好看
\begin{tikzpicture}[remember picture,overlay]
\node[inner sep=5cm,outer sep=5cm,opacity=0.1]%0.1是透明度，调节到正常放映
  at (current page.center)
  {\includegraphics[width=6cm,height=6cm]{logo.jpg}};
\end{tikzpicture}
}
\setbeamertemplate{bibliography item}[text]
%参考文献正常编号


\mode<presentation> {
	\usetheme{Madrid}
 %可选代替Madrid：
 % AnnArbor Antibes Bergen Berkeley Berlin Boadilla CambridgeUS Copenhagen Darmstadt 
 % Dresden EastLansing Frankfurt Goettingen Hannover Ilmenau JuanLesPins Luebeck  
 % Malmoe Marburg Montpellier PaloAlto Pittsburgh Rochester Singapore Szeged Warsaw


	\usecolortheme{spruce}
    \insertlogo
    \setbeamercolor{titlelike}{fg=paris}
    \setbeamercolor{structure}{fg=paris}
	%修改成巴院LOGO颜色


    \setbeamertemplate{navigation symbols}{}
    %要删除所有幻灯片底部的导航符号，请取消注释此行
	
	%\setbeamertemplate｛footline｝
	%若要删除所有幻灯片中的页脚，请取消注释此行，可以避免答辩老师一眼看出没有认真做
	
	%\setbeamertemplate｛footline｝[页码]
	%若要用简单的幻灯片计数替换所有幻灯片中的页脚，请取消注释此行，嗯嗯
}

\usepackage{graphicx} % 允许包含图像
\usepackage{booktabs} % 允许在表中使用\toprule、\ midrule和\ bottomrule
\usepackage{beamerfoils}
\usepackage[UTF8,noindent]{ctexcap}  % 使用中文输入及显示
\usepackage[bookmarks=true]{hyperref}
\MyLogo{\vspace*{8.5cm}\includegraphics[height=0.7cm]{BUCT-heng2.jpg}\hspace{0.1cm}\includegraphics[height=0.7cm]{zuhe.jpg}} 
%\setbeamercolor{block title}{fg = white,bg = paris}
%-----------------------------------
%	以下为正文
%-----------------------------------

\title[开题答辩]{{工程师学院Beamer模板研究} }
% 简短标题显示在每张幻灯片的底部，完整标题仅在标题页上
\subtitle{{以一个绿色样式为例}}
\author{Sihidyom} % Your name
\institute[] % 您的机构将出现在每张幻灯片的底部，可能是节省空间的简写
{	
	工生2077 2077100114\\
	北京化工大学巴黎居里工程师学院 \\ % 您所在的机构
    \vspace{0.5cm}

	指导教师：AAA$^1$、BBB$^2$\\
	$^1$ 北京化工大学巴黎居里工程师学院\\
	$^2$ CC大学DD学院
	%\medskip
	%\textit{zyl@buct.edu.cn} 
    % email address
}
%\date{\today} % 日期，可以更改为自定义日期

\begin{document}
	
	\begin{frame}
		\titlepage % 将标题页打印为第一张幻灯片
	\end{frame}
	\LogoOff % 以下不显示logo
	\begin{frame}
        \Background
		\frametitle{目录} % 目录幻灯片，注释此块以将其删除
		\tableofcontents % 在整个演示过程中，如果您选择使用\ section｛｝和\ submission｛}命令，这些命令将自动打印在此幻灯片上，作为演示的概述
	\end{frame}
	
	%-----------------------------------
	%	开始创建
	%-----------------------------------
	
\section{研究背景} 
\subsection{consectetur adipiscing}
		\begin{frame}\Background
			\frametitle{Lorem ipsum dolor sit amet}
			\begin{columns}
				\column{0.4\textwidth}
				\begin{figure}
					\includegraphics[width=1\textwidth]{zuhe.jpg}
					\caption{如彭奇和瓦特曼公业\cite{1}}
				\end{figure}

				\column{0.6\textwidth}
				Lorem ipsum dolor sit amet, consectetur adipiscing elit. Integer dolor elit, auctor a dictum ut, lobortis non ex. Proin lectus dolor, viverra at eros eu, egestas eleifend urna. Mauris risus nisl, auctor quis viverra at, tempus vel risus. 
			\end{columns}
			
		\end{frame}


        \subsection{Phasellus tincidunt}
		\begin{frame}\Background
			\frametitle{Phasellus tincidunt id nibh ac dictum.}
			\vspace{-0.5cm}
			\begin{block}{ Duis tempus volutpat neque eget vehicula.}
						
				\begin{columns}
					\column{0.5\textwidth}
					\begin{figure}
						\includegraphics[width=0.6\textwidth]{BUCT-heng2.jpg}
						\caption{Curabitur tempus }
					\end{figure}

				\column{0.5\textwidth}
				\begin{itemize}
					\item Integer sollicitudin turpis ultrices. 
					\item Etiam quis velit metus. 
				\end{itemize}
				\end{columns}
			\end{block}

			\begin{block}{Nulla malesuada}
				\begin{columns}
					\column{0.5\textwidth}
					\begin{figure}
						\includegraphics[width=0.6\textwidth]{zuhe.jpg}
						\caption{Nulla malesuada magna scelerisque mauris condimentum fermentum. \cite{Suspendisse}}
					\end{figure}
					\column{0.5\textwidth}	
					\begin{itemize}
					\item Morbi placerat dolor lacus, in tempor lacus porta in. 
					\item Aliquam leo eros, volutpat et lectus vitae.
				\end{itemize}
				\end{columns}
			\end{block}
		\end{frame}


        \subsection{Suspendisse}
		\begin{frame}\Background
			\frametitle{Suspendisse ullamcorper accumsan lectus in malesuada. Morbi et luctus turpis.}
			\begin{itemize}
				\item Aliquam non erat eu tellus elementum convallis.
				\item Aliquam non erat eu tellus elementum convallis.
				\item Aliquam non erat eu tellus elementum convallis.
			\end{itemize}
			\begin{figure}
				\includegraphics[width=0.9\textwidth]{BUCT-heng.jpg}
				\caption{Neque\cite{Neque}}
			\end{figure}
			
		\end{frame}
		\begin{frame}\Background
			\frametitle{fringilla}
			\begin{columns}
				\column{0.45\textwidth}
				\begin{block}{目前的主流研究和方向}
					\begin{itemize}
						\item 去讲哎哟
						\item 康纳马拉尽
						\item 寒冷中在海
						\item fringilla risus
						\item 普汉换
						\item 大概约略
					\end{itemize}
				\end{block}
				
				\column{0.45\textwidth}
				\begin{block}{前人的研究和不足}
					\begin{itemize}
						\item item
						\item itemize
                            \item emurate
					\end{itemize}
				\end{block}
				\begin{block}{本课题主要创新点}
					\begin{itemize}
						\item 
						\item 
						\item 
						\item 
					\end{itemize}
				\end{block}

			\end{columns}
			
		\end{frame}
 \section{研究计划}
    \subsection{实验安排}
    \subsubsection{a.}
    \subsubsection{b.}
    \subsubsection{c.}
    \subsection{结构表征}
    \subsubsection{核磁共振氢谱}
    \subsubsection{傅里叶红外分析}
    \subsubsection{X射线衍射}
    \subsection{数据处理}
    \subsection{时间进度}
    \subsection{预期展望}
 \section{参考文献}
	\begin{frame}[allowframebreaks]\Background
		\frametitle{参考文献}
		\begin{thebibliography}{99}
			\bibitem[1]{1} In posuere vitae urna non consequat. Aliquam dapibus, enim ultrices.
			\bibitem[2]{Suspendisse} Suspendisse ullamcorper accumsan lectus in malesuada. Morbi et luctus turpis.
			\bibitem[3]{Neque} Neque porro quisquam est qui dolorem ipsum quia dolor sit amet, consectetur, adipisci velit.
                \bibitem[4]{squamest} allllllllllllowframebreaks
			\end{thebibliography}
		
	\end{frame}

	\LogoOn
	\begin{frame}\Background
		\Huge{\centerline{感谢各位老师的批评指正!}}
		\centering{{\normalsize }}
	\end{frame}
	
\end{document} 

\end{document}